\chapter{Seeding}
\label{ch:seeding}

Seeding turns a read's raw sketch (a flat list of $k$-mers, with repeats) into an ordered list of
\emph{seeds} --- one entry per distinct hash --- ready to be consumed budget-first against the
index. This chapter covers \code{unique\_elements\_with\_info} (grouping) and the derivation of
the seed budget $S$ used by the caller.

\section{Grouping: \code{unique\_elements\_with\_info}}
\label{sec:seeding-grouping}

\paragraph{Input.} \code{p}: the read's sketch, \code{Vec<Kmer>} (Chapter~\ref{ch:sketching}),
mutable (it is sorted in place). \emph{Implicit input:} the reference index, queried once per
distinct hash.

\paragraph{Output.} \code{Vec<Seed>} (\S\ref{sec:type-seed}), one per distinct hash value present
in \code{p}, \textbf{sorted ascending by \code{hits\_in\_t}} (rarest-in-reference first). Also
has the side effect of updating the \code{kmers\_notmatched} diagnostic counter
(\S\ref{sec:defect-notmatched} documents a bug in how the reference computes this specific
counter --- it does not affect any other output).

\begin{algobox}{title={\code{unique\_elements\_with\_info}(p)}}
\begin{verbatim}
# Step 1: group equal hashes together, and within a group, order by
# *decreasing* read position (this ordering becomes each seed's pmatches
# list, which the bucket_LCS metric consumes directly, see 10.4).
sort p by (h ascending, r descending)

# Step 2: walk the sorted list, closing out a group whenever the hash changes.
p_unique := []
strike := 0                 # occurrences seen in the current group so far
matches := []                # positions seen in the current group so far
for pos in 0 .. len(p)-1:
    strike += 1
    matches.append(p[pos].r)
    if pos == len(p)-1 or p[pos].h != p[pos+1].h:
        hits_in_t := index.count(p[pos].h)                # Section 6.4
        p_unique.append(Seed{
            kmer: p[pos], hits_in_t: hits_in_t,
            occs_in_p: strike, seed_num: len(p_unique),
            pmatches: copy(matches),
        })
        strike := 0
        matches := []

# Step 3: seeds are consumed rarest-reference-hit-first.
sort p_unique by hits_in_t ascending
return p_unique
\end{verbatim}
\end{algobox}

\paragraph{Why sort by \code{hits\_in\_t} ascending.} A $k$-mer that occurs rarely in the
reference is far more discriminative than a common one: matching it narrows the candidate set of
buckets dramatically, whereas matching a $k$-mer with thousands of reference occurrences adds
thousands of (mostly irrelevant) candidate hits for very little narrowing benefit. Consuming
seeds rarest-first under a fixed budget $S$ (\S\ref{sec:seed-budget}) maximizes the discriminative
value extracted from that budget.

\paragraph{Why sort the group by decreasing position.} \code{Seed.pmatches} ends up populated in
the same order the group was scanned, i.e.\ decreasing read position --- this specific order is
required by \code{lcs\_get\_lis} (\S\ref{sec:lcs-implementation}), which needs \code{pmatches} in
a defined order to reconstruct multiplicities correctly across repeated seeds when computing a
longest increasing subsequence.

\subsection{Reference C++ implementation}

\begin{lstlisting}[style=cppstyle,caption={\code{unique\_elements\_with\_info}, \texttt{src/shmap.h}}]
Seeds unique_elements_with_info(sketch_t& p) {
    sort(p.begin(), p.end(), [](const Kmer &a, const Kmer &b) {
        if (a.h != b.h) return a.h < b.h;
        return a.r > b.r;                 // reverse order of inclusion; needed for LCS
    });

    Seeds p_unique;
    qpos_t strike = 0;
    vector<qpos_t> matches;
    for (size_t ppos = 0; ppos < p.size(); ++ppos) {
        ++strike;
        matches.push_back(p[ppos].r);
        if (ppos == p.size()-1 || p[ppos].h != p[ppos+1].h) {
            rpos_t hits_in_t = tidx.count(p[ppos].h);
            p_unique.emplace_back(p[ppos], hits_in_t, strike, p_unique.size(), matches);
            strike = 0;
            matches.clear();
        }
    }

    sort(p_unique.begin(), p_unique.end(), [](const Seed &a, const Seed &b) {
        return a.hits_in_T < b.hits_in_T;
    });
    return p_unique;
}
\end{lstlisting}

\begin{bugbox}
\label{sec:defect-notmatched}
The reference C++ additionally maintains a \code{nonzero} counter meant to track how many
\emph{read positions} (counting multiplicity) belong to seeds that do have at least one reference
hit, so that the diagnostic \code{kmers\_notmatched} stat can report ``sketch size minus
matched positions''. The increment is written as \code{if (hits\_in\_t > 0) nonzero += strike;}
placed \emph{after} \code{strike} has already been reset to $0$ for the next group two lines
above it --- so it always adds $0$, and \code{kmers\_notmatched} always reports the entire sketch
size regardless of how many $k$-mers actually matched. This affects \emph{only} a cosmetic,
end-of-run diagnostic percentage; it has no effect on any mapping decision, score, or PAF field.
A from-scratch implementation should compute the increment \emph{before} resetting \code{strike}
(or reorder the two statements), which is the obviously-intended behavior.
\end{bugbox}

\section{The seed budget $S$}
\label{sec:seed-budget}

Not every seed is necessarily consumed --- seeding stops once a budget $S$ of seed
\emph{occurrences} (not distinct seeds; a repeated $k$-mer with \code{occs\_in\_p}$=3$ counts as
$3$ against the budget when consumed, all at once, \S\ref{sec:match-seeds}) has been spent. $S$ is
derived once per read from the acceptance threshold $\theta$ (CLI \code{-t}) and the
minimum-score-gap parameter (CLI \code{-d}, used here as a safety margin, not in its other role
computing the second-best threshold in Chapter~\ref{ch:end-to-end}):
\[
  \theta_2 = \theta - \mathit{min\_diff}, \qquad
  S = \left\lfloor (1-\theta_2)\cdot m \right\rfloor + 1
\]
where $m$ is the read's full sketch size. \textbf{Interpretation:} $S$ is exactly the smallest
seed-occurrence budget for which it is still \emph{possible}, in principle, for a bucket to reach
score $\theta_2$ even if literally every consumed seed occurrence that has a reference hit at all
ends up matching the same bucket: the seed-heuristic bound (\S\ref{sec:seed-heuristic-theory})
for a bucket that has consumed all $S$ seed occurrences and matched every one of them is exactly
$h_{\mathrm{seed}} = 1 - 0/m$, and the bound only starts to allow room for \emph{misses} once more
than $(1-\theta_2)m$ occurrences have been consumed --- so $S$ is the minimum budget that gives the
pruning bound any chance of staying above $\theta_2$ for a maximally-good bucket, with $+1$
guaranteeing at least one seed is always consumed even when $\theta_2 \to 1$.

\begin{notebox}
This is a heuristic sizing choice, not a hard correctness requirement in either direction: seeds
\emph{beyond} $S$ might still have found valuable matches (this is a real, accepted
recall/speed trade-off, and is exactly the knob CLI flag \code{-t} tunes indirectly through
$\theta_2$); a from-scratch implementation experimenting with alternative sizing should keep in
mind that $S$ only bounds how many seed \emph{occurrences} are matched against the index at all
--- it does not, by itself, bound the pruning behavior described in Chapter~\ref{ch:pruning},
which operates independently per bucket using whatever prefix of seeds \emph{has} been consumed
so far.
\end{notebox}

\section{\code{more\_seeds\_if\_cheap}: opportunistic budget extension (dead code)}

\begin{algobox}{title={\code{more\_seeds\_if\_cheap}(S, seeds)}}
\begin{verbatim}
original_S := S
while S < len(seeds) and seeds[S].hits_in_t <= 1:
    S += 1
return S
\end{verbatim}
\end{algobox}

This extends the budget past $S$ for free whenever the next seed(s) are (almost) unique in the
reference (\code{hits\_in\_t} $\le 1$): such seeds are essentially free to consume (at most one
reference hit to check) and maximally discriminative, so there's no real cost/benefit trade-off in
including them.

\begin{bugbox}
This function is fully implemented and unit-testable, but its only call site in the reference C++
\code{map\_read} is commented out (\code{//S = more\_seeds\_if\_cheap(S, p\_unique);}) --- it is
never actually invoked by the shipped tool. A from-scratch implementation may freely choose to
wire it in (it is a strict improvement in recall at negligible extra cost) or leave it unused to
match the reference binary's actual observed behavior exactly.
\end{bugbox}
